\documentclass{zufe}
\usepackage{setspace} % 用于设置行间距
%首航缩进
\usepackage{indentfirst}
%各种下划线和url
\usepackage{hyperref}
\usepackage{ulem}
% 图表设置
\usepackage{graphicx}
\usepackage{caption}
\usepackage{appendix}
\usepackage{lipsum}
%数学相关包
\usepackage{algorithm}
\usepackage[noend]{algpseudocode} % 使用 algpseudocode 包提供的接口
% 自定义跨页算法环境
\makeatletter
\newenvironment{breakablealgorithm}
{% \begin{breakablealgorithm}
		\begin{center}
			\refstepcounter{algorithm}% New algorithm
			\hrule height.8pt depth0pt \kern2pt% \@fs@pre for \@fs@ruled
			\renewcommand{\caption}[2][\relax]{% Make a new \caption
				{\raggedright\textbf{\textbf{Algorithm}~\thealgorithm} ##2\par}%
				\ifx\relax##1\relax % #1 is \relax
				\addcontentsline{loa}{algorithm}{\protect\numberline{\thealgorithm}##2}%
				\else % #1 is not \relax
				\addcontentsline{loa}{algorithm}{\protect\numberline{\thealgorithm}##1}%
				\fi
				\kern2pt\hrule\kern2pt
			}
		}{% \end{breakablealgorithm}
		\kern2pt\hrule\relax %\@fs@post% for \@fs@ruled
	\end{center}
}
\makeatother
%====================Tiks ====================
\usepackage{tikz}
\usetikzlibrary{shapes.geometric, arrows.meta, positioning}
\tikzstyle{process} = [rectangle, minimum width=3cm, minimum height=1cm, text centered, draw=black, fill=blue!10]
\tikzstyle{arrow} = [thick, ->, >=stealth]
%====================Tiks ====================

\usepackage{amsmath} % 支持数学符号
\usepackage{amssymb}
\usepackage{algorithmicx}
%处理数学公式中和黑斜体
\usepackage{latexsym,bm}
%定义书写公式环境
\usepackage{cases}
%一些特殊符号例如：※√
\usepackage{pifont}

% 数学命令
\makeatletter
\newcommand\dif{%  % 微分符号
	\mathop{}\!%
	\ifthu@math@style@TeX
	d%
	\else
	\mathrm{d}%
	\fi
}
\makeatother
% 表格中支持跨行
\usepackage{multirow}

% 固定宽度的表格。
\usepackage{tabularx}

% 跨页表格
\usepackage{longtable}
%给字母加粗
\usepackage[T1]{fontenc}

\usepackage{amsthm}
% 量和单位
\usepackage{siunitx}

\newcommand{\formatmain}{
	\setlength{\parskip}{0em}
	\renewcommand{\baselinestretch}{1.53}
	\interfootnotelinepenalty=10000}
%三线表
\usepackage{booktabs}

%页眉页脚
\usepackage{fancyhdr}
%三线表
\usepackage{threeparttable}
% 中英文摘要和关键字
%这里标题的设置是黑体 三号 居中
%新罗马字体
\newcommand{\romann}{\CJKfontspec{TIMES .TTF}\selectfont}
\newtheorem{theorem}{定理}

%重命名lst名字
\renewcommand{\lstlistingname}{代码}
%参考文献生成规则
\usepackage[
backend=biber,
style=gb7714-2015,
gbalign=gb7714-2015,
gbnamefmt=lowercase,
gbpub=false,
doi=false,
url=false,
eprint=false,
isbn=false,
]{biblatex}

\hypersetup{
	colorlinks=true,
	linkcolor=black,
	citecolor=black,
	urlcolor=black
}


\addbibresource{ref/refs.bib}


% % 使用中文编译 + XeLatex + Biber
% \usepackage[fontset=ubuntu]{ctex}

% 基本信息---------------------------------------------------------------
% 在这里填写你的论文中英文文题目
\newcommand{\thesisname}{基于latex的zufe论文模板}
\newcommand{\subtitle}{副标题示例} %如果没有副标题，注掉
\newcommand{\enthesisname}{English Main Title Example}  %英文主标题，如果不需要英文摘要就注掉，但是同时摘要里面的英文摘要也就没了（英文标题没了应该没人会写英文摘要吧）
\newcommand{\ensubtitle}{English Subtitle Example}
\newcommand{\studentname}{学生姓名}
\newcommand{\studentid}{学号示例}
\newcommand{\major}{专业示例}
\newcommand{\advisor}{指导教师示例}
\newcommand{\deptName}{学院名称}
\newcommand{\className}{班级名称}
\newcommand{\thesisdate}{2024年6月}

\begin{document}
	
	% 生成封面
	% 自定义下划线命令
\newcommand\dunderline[3][-1pt]{{%
		\setbox0=\hbox{#3}%
		\ooalign{\copy0\cr\rule[\dimexpr#1-#2\relax]{\wd0}{#2}}}}

% 生成封面
\begin{titlepage}
	\makeatletter
	%是否有校外指导老师，有的话就运行12行，否则运行15行，下面42行处的判断同理
	\@ifundefined{externalMentorName}{
		% 校内毕设封面顶部间距
		\vspace*{-20mm}
	}{
		% 校外毕设封面顶部间距
		\vspace*{-15mm}
	}
	\centering
	
	\includegraphics[width=3cm]{figures/logo.png}\\
	\includegraphics[width=5cm, height=1cm]{figures/name.png}
	\vspace*{10mm}
	\begin{center}
  		\zihao{1}\textbf{\ziju{0.12}\songti{本\hspace{5mm}科\hspace{5mm}生\hspace{5mm}毕\hspace{5mm}业\hspace{5mm}论\hspace{5mm}文（设计）}}
		
		\vspace{20mm}
		
		\xiaoer\textmd{\heiti{题目：}\thesisname}
		
		\vspace{5mm}
		%判断是否有副标题
		\ifx\subtitle\undefined\else
		%如果有副标题
		\begin{spacing}{1.2}
			\sanhao\selectfont{\textmd{\kaishu{-----}\subtitle}}
		\end{spacing}
		\fi
		\vspace{30mm}
		
		\flushleft
		\makeatletter
		\@ifundefined{externalMentorName}
		{
			\makeatother
			% 生成校内毕设封面字段
			\begin{spacing}{1.8}
				\hspace{27mm}\songti\zihao{3}\selectfont{学生姓名：\dunderline[-10pt]{1pt}{\makebox[78mm][c]{\studentname}}}
				
				\hspace{27mm}\songti\zihao{3}\selectfont{学\hspace{11mm}号：\dunderline[-10pt]{1pt}{\makebox[78mm][c]{\studentid}}}
				
				\hspace{27mm}\songti\zihao{3}\selectfont{指导教师：\dunderline[-10pt]{1pt}{\makebox[78mm][c]{\advisor}}}
				
				\hspace{27mm}\songti\zihao{3}\selectfont{所在学院：\dunderline[-10pt]{1pt}{\makebox[78mm][c]{\deptName}}}
				
				\hspace{27mm}\songti\zihao{3}\selectfont{专业名称：\dunderline[-10pt]{1pt}{\makebox[78mm][c]{\major}}}
				
				\hspace{27mm}\songti\zihao{3}\selectfont{班\hspace{11mm}级：\dunderline[-10pt]{1pt}{\makebox[78mm][c]{\className}}}
			\end{spacing}
		}
		{
			\makeatother
			% 生成校外毕设封面字段
			\begin{spacing}{1.8}
				\hspace{19.4mm}\songti\zihao{3}\selectfont{学\hspace{11mm}院：\dunderline[-10pt]{1pt}{\makebox[77.4mm][c]{\deptName}}}
				
				\hspace{19.4mm}\songti\zihao{3}\selectfont{专\hspace{11mm}业：\dunderline[-10pt]{1pt}{\makebox[77.4mm][c]{\major}}}
				
				\hspace{19.4mm}\songti\zihao{3}\selectfont{学生姓名：\dunderline[-10pt]{1pt}{\makebox[77.4mm][c]{\studentname}}}
				
				\hspace{19.4mm}\songti\zihao{3}\selectfont{学\hspace{11mm}号：\dunderline[-10pt]{1pt}{\makebox[77.4mm][c]{\studentid}}}
				
				\hspace{19.4mm}\songti\zihao{3}\selectfont{指导教师：\dunderline[-10pt]{1pt}{\makebox[77.4mm][c]{\advisor}}}
			\end{spacing}
		}

		
		\vspace{25mm}
		\centering
		\zihao{4}\ziju{0.5}{\thesisdate}
	\end{center}
\end{titlepage}


	\includepdf[pages=-]{figures/copyright(2).pdf}

	\frontmatter
	% 摘要（中英文）：根据自身论文，修改摘要的Tex文件
	
% 下面是中文摘要 3-19行不用管
\setstretch{2}
\vspace*{2em}

\begin{center}
	\zihao{3}\textbf{ABSTRACT}
	\vspace{2mm}
	\ifx\subtitle\undefined\else
	%如果有副标题
	\begin{spacing}{1.2}
		\sanhao\selectfont{\textmd{\kaishu{-----}\subtitle}}
	\end{spacing}
	\fi
\end{center}
%题目后空一行
\vspace{1em}
\setstretch{1.5}

\setlength{\parskip}{0em}

% 英文摘要正文从这里开始--------------------------------------------------
As economic globalization continues to advance in depth, the translation of academic literature on economics has increasingly demonstrated its significance. This report selects and translates the 2015 Nobel Prize laureate’s academic literature into Chinese. Based on the research of the selected text and its Chinese translation by combining the case analysis with the comparative analysis, the report attempts to focus on what translation strategies are more effective in terms of realizing the ideational, interpersonal, and textual functions of the source text. The selected materials of this translation practice report do not have any existing translation versions for reference and the present author uses computer-aided technology to conduct a data analysis of the full original text before translation so as to learn about the features of the original.
\par
This report is guided by the three meta-functions of Halliday’s systemic functional linguistics. Through the analysis, it is concluded that annotation of terminologies, division of long sentences, and conversion of passive voice can achieve the ideational functions of the translation; addition of personal pronouns, literal translation of modal verbs and substitution of mood adjuncts can realize the interpersonal function; omission of connectives and conjunctive adjuncts, reproduction of demonstrative pronouns and adverbs, and reorganization of sentence order and paragraphs can achieve the textual function. The practicability of these strategies has been confirmed by concrete examples in this report. For the translation of academic literature on economics, comprehension and expression are pivotal. The translator should be familiar with the characteristics and style of the text, proficient in translation skills with the help of computer-aided tools, and flexible in dealing with rendering difficulties. This report is also intended to give some implications to both the translation practice of and theoretical application to the similar texts in the future.
\par

\vspace{1em}
\textbf{\heiti\xiaosi{Key words:}}
%这里写你的关键词------------------------------------------------------
academic literature on economics; three meta-functions of language; translation strategies

\newpage
\setstretch{2}
\vspace*{2em}
%如果没有英文标题，跳过
%下面是英文摘要 37-52行不用管
\begin{spacing}{0.95}
	\centering
	\sanhao\textbf{\heiti{摘要}}
	\vspace{2mm}
	\ifx\subtitle\undefined\else
	%如果有副标题
	\begin{spacing}{1.2}
		\sihao\selectfont{\textmd{\kaishu{-----}\rmfamily\textbf{\ensubtitle}}}
	\end{spacing}
	\fi
\end{spacing}
\vspace{1em}
\setstretch{1.5}
\setlength{\parskip}{0em}

% 中文摘要正文从这里开始--------------------------------------------------
随着经济全球化不断纵深推进，经济学类学术文献的翻译愈发显示出其重要性。本报告选择2015年诺贝尔经济学奖获得者的学术讲演稿进行汉译，基于对所选文本及其汉译文的研究，结合个案分析法与对比分析法，试图重点考察经济学类学术文献的翻译过程中采用何种策略可以在译文中更为有效地实现概念功能、人际功能、语篇功能等三大语言纯理功能。本翻译实践报告的所选材料无任何现有译文可供参考，笔者为了解原文特点，在译前借助计算机辅助技术对全文进行了数据分析。
\par
本报告以韩礼德系统功能语言学理论的三大语言纯理功能为指导，通过分析，结论如下：对术语加注、拆分长句和转换被动语态能够实现译文的概念功能；增译人称代词、直译情态动词和替换语气附加语可以实现译文的人际功能；省却连词和连接附加语、再现指示代词和指示副词，以及重组句序和段落能够实现译文的语篇功能，这些策略的实用性已由本报告的翻译实例得到印证。对于经济学类学术文献的翻译，理解能力与表达方式至关重要，译者当熟悉文体特征，精于翻译技巧并辅以机辅工具，灵活应对翻译难点。本报告对今后类似文本的翻译实践及理论应用也具有一定的启示意义。
\par

\vspace{1em}
\textbf{\rmfamily\xiaosi{关键词：}}
%这里写你的英文关键词----------------------------------------------------
经济学类学术文献；三大语言纯理功能；翻译策略



	\newpage
	%这个是目录页，不用修改
	% 论文目录
% 没有特殊需要不用修改
% 目录
\addtocontents{toc}{\protect\thispagestyle{ZUFEThesis}}
\tableofcontents
\clearpage
	
	% 主要内容部分
	\mainmatter
	\formatmain
	
	%这里是引言
	% 引言
\chapter*{Introduction}
\addcontentsline{toc}{chapter}{Introduction}
Economic development in the new era has played a more prominent role than ever in personal life, social change, country prospect as well as international stage. Whether the ultimate aim is to be in tandem with or outperform other economies, nearly every single country and region strives to engage in economic activities and discussions worldwide, therefore, economic texts, especially academic literature on economics have witnessed a boom of multilingual translation and cross-border dissemination, effectively and efficiently bridging language gaps and exchanging economic philosophies.

...


	
	%第一章
	
\chapter{第一章}
\section{论文主要部分的写法}

研究生学位论文撰写，除表达形式上需要符合一定的格式要求外，内容方面上也要遵循一些共性原则。

通常研究生学位论文只能有一个主题（不能是几块工作拼凑在一起），该主题应针对某学科领域中的一个具体问题展开深入、系统的研究，并得出有价值的研究结论。
学位论文的研究主题切忌过大，例如，“中国国有企业改制问题研究”这样的研究主题过大，因为“国企改制”涉及的问题范围太广，很难在一本研究生学位论文中完全研究透彻。



\subsection{论文的语言及表述}

除国际研究生外，学位论文一律须用汉语书写。
学位论文应当用规范汉字进行撰写，除古汉语研究中涉及的古文字和参考文献中引用的外文文献之外，均采用简体汉字撰写。

国际研究生一般应以中文或英文书写学位论文，格式要求同上。
论文须用中文封面。

研究生学位论文是学术作品，因此其表述要严谨简明，重点突出，专业常识应简写或不写，做到立论正确、数据可靠、说明透彻、推理严谨、文字凝练、层次分明，避免使用文学性质的或带感情色彩的非学术性语言。

论文中如出现一个非通用性的新名词、新术语或新概念，需随即解释清楚。



\subsection{论文题目的写法}

论文题目应简明扼要地反映论文工作的主要内容，力求精炼、准确，切忌笼统。
论文题目是对研究对象的准确、具体描述，一般要在一定程度上体现研究结论，因此，论文题目不仅应告诉读者这本论文研究了什么问题，更要告诉读者这个研究得出的结论。
例如：“在事实与虚构之间：梅乐、卡彭特、沃尔夫的新闻观”就比“三个美国作家的新闻观研究”更专业、更准确。



\subsection{摘要的写法}

论文摘要是对论文研究内容的高度概括，应具有独立性和自含性，即应是 一篇简短但意义完整的文章。
通过阅读论文摘要，读者应该能够对论文的研究 方法及结论有一个整体性的了解，因此摘要的写法应力求精确简明。
论文摘要 应包括对问题及研究目的的描述、对使用的方法和研究过程进行的简要介绍、 对研究结论的高度凝练等，重点是结果和结论。

论文摘要切忌写成全文的提纲，尤其要避免“第 1 章……；第 2 章……；……”这样的陈述方式。



\subsection{引言的写法}

一篇学位论文的引言大致包含如下几个部分：
1、问题的提出；
2、选题背 景及意义；
3、文献综述；
4、研究方法；
5、论文结构安排。
\begin{itemize}
	\item 问题的提出：要清晰地阐述所要研究的问题“是什么”。
	\footnote{选题时切记要有“问题意识”，不要选不是问题的问题来研究。}
	\item 选题背景及意义：论述清楚为什么选择这个题目来研究，即阐述该研究对学科发展的贡献、对国计民生的理论与现实意义等。
	\item 文献综述：对本研究主题范围内的文献进行详尽的综合述评，“述”的同时一定要有“评”，指出现有研究状态，仍存在哪些尚待解决的问题，讲出自己的研究有哪些探索性内容。
	\item 研究方法：讲清论文所使用的学术研究方法。
	\item 论文结构安排：介绍本论文的写作结构安排。
\end{itemize}



\subsection{正文的写法}

本部分是论文作者的研究内容，不能将他人研究成果不加区分地掺和进来。
已经在引言的文献综述部分讲过的内容，这里不需要再重复。
各章之间要存在有机联系，符合逻辑顺序。



\subsection{结论的写法}

结论是对论文主要研究结果、论点的提炼与概括，应精炼、准确、完整，使读者看后能全面了解论文的意义、目的和工作内容。
结论是最终的、总体的结论，不是正文各章小结的简单重复。
结论应包括论文的核心观点，主要阐述作者的创造性工作及所取得的研究成果在本领域中的地位、作用和意义，交代研究工作的局限，提出未来工作的意见或建议。
同时，要严格区分自己取得的成果与指导教师及他人的学术成果。

在评价自己的研究工作成果时，要实事求是，除非有足够的证据表明自己的研究是“首次”、“领先”、“填补空白”的，否则应避免使用这些或类似词语。

	
	%下面自己按照自己的论文增加章节……
	% !TeX root = ../zufe_report.tex
\setlength{\baselineskip}{20pt}
\section{表格示例}
\subsection{表格}

表应具有自明性。为使表格简洁易读，尽可能采用三线表，如表~\ref{tab:three-line}。
三条线可以使用 booktabs 宏包提供的命令生成。

\begin{table}[htbp]
	\centering
	\caption{三线表示例}
	\begin{tabular}{ll}
		\toprule
		文件名          & 描述                         \\
		\midrule
		zufereport.dtx   & 模板的源文件，包括文档和注释 \\
		zufereport.cls   & 模板文件                     \\
		zufereport-*.bst & BibTeX 参考文献表样式文件    \\
		\bottomrule
	\end{tabular}
	\label{tab:three-line}
\end{table}

表格如果有附注，尤其是需要在表格中进行标注时，可以使用 threeparttable 宏包。
研究生要求使用英文小写字母 a、b、c……顺序编号，本科生使用圈码 ①、②、③……编号。

\begin{table}[htbp]
	\centering
	\begin{threeparttable}[c]
		\caption{带附注的表格示例}
		\label{tab:three-part-table}
		\begin{tabular}{ll}
			\toprule
			文件名                 & 描述                         \\
			\midrule
			zufereport.dtx\tnote{a} & 模板的源文件，包括文档和注释 \\
			zufereport.cls\tnote{b} & 模板文件                     \\
			zufereport-*.bst        & BibTeX 参考文献表样式文件    \\
			\bottomrule
		\end{tabular}
		\begin{tablenotes}
			\item [a] 可以通过 xelatex 编译生成模板的使用说明文档；
			使用 xetex 编译 \file{zufereport.ins} 时则会从 \file{.dtx} 中去除掉文档和注释，得到精简的 \file{.cls} 文件。
			\item [b] 更新模板时，一定要记得编译生成 \file{.cls} 文件，否则编译论文时载入的依然是旧版的模板。
		\end{tablenotes}
	\end{threeparttable}
\end{table}

如某个表需要转页接排，可以使用 \usepackage{longtable} 宏包，需要在随后的各页上重复表的编号。
编号后跟表题（可省略）和“（续）”，置于表上方。续表均应重复表头。

\singlespacing
\begin{longtable}{cccc}
	\caption{跨页长表格的表题}
	\label{tab:longtable} \\
	\toprule
	表头 1 & 表头 2 & 表头 3 & 表头 4 \\
	\midrule
	\endfirsthead
	\caption*{续表~\thetable\quad 跨页长表格的表题} \\
	\toprule
	表头 1 & 表头 2 & 表头 3 & 表头 4 \\
	\midrule
	\endhead
	\bottomrule
	\endfoot
	Row 1  & & & \\
	Row 2  & & & \\
	Row 3  & & & \\
	Row 4  & & & \\
	Row 5  & & & \\
	Row 6  & & & \\
	Row 7  & & & \\
	Row 8  & & & \\
	Row 9  & & & \\
	Row 10 & & & \\
\end{longtable}


\subsection{算法}

算法环境可以使用 algorithms 或者 algorithm2e 宏包。

\renewcommand{\algorithmicrequire}{\textbf{输入：}\unskip}
\renewcommand{\algorithmicensure}{\textbf{输出：}\unskip}

\begin{algorithm}
	\caption{Calculate $y = x^n$}
	\label{alg1}
	\small
	\begin{algorithmic}
		\REQUIRE $n \geq 0$
		\ENSURE $y = x^n$
		
		\STATE $y \leftarrow 1$
		\STATE $X \leftarrow x$
		\STATE $N \leftarrow n$
		
		\WHILE{$N \neq 0$}
		\IF{$N$ is even}
		\STATE $X \leftarrow X \times X$
		\STATE $N \leftarrow N / 2$
		\ELSE[$N$ is odd]
		\STATE $y \leftarrow y \times X$
		\STATE $N \leftarrow N - 1$
		\ENDIF
		\ENDWHILE
	\end{algorithmic}
\end{algorithm}

	
	%数学公式以及表格和算法
	% !TeX root = ../zufe_report.tex
\setlength{\baselineskip}{20pt}

\section{数学公式部分}
\subsection{数学符号和公式}
\subsubsection{标准}
中文论文的数学符号默认遵循 GB/T 3102.11—1993《物理科学和技术中使用的数学符号》
\footnote{原 GB 3102.11—1993，自 2017 年 3 月 23 日起，该标准转为推荐性标准。}。
该标准参照采纳 ISO 31-11:1992 \footnote{目前已更新为 ISO 80000-2:2019。}，
但是与 \TeX{} 默认的美国数学学会（AMS）的符号习惯有所区别。
具体地来说主要有以下差异：
\begin{enumerate}
	\item 
	大写希腊字母默认`为斜体，如
	\begin{equation*}
		\varGamma \varDelta \varTheta \varLambda \varXi \varPi \varSigma \varUpsilon \varPhi \varPsi \varOmega.
	\end{equation*}
	\item 小于等于号和大于等于号使用倾斜的字形 $\le$、$\ge$。
	\item 积分号使用正体，比如 $\int$、$\oint$。
	\item
	偏微分符号 $\partial$ 使用正体。
	\item
	省略号按照中文的习惯固定居中，比如
	\begin{equation*}
		1, 2, \cdots, n \quad 1 + 2 + \cdots + n.
	\end{equation*}
	\item
	实部 $\operatorname{Re}$ 和虚部 $\operatorname{Im}$ 的字体使用罗马体。
\end{enumerate}




\subsubsection{TeX{}基础操作}
\begin{enumerate}
	\item TeX{}中的行内公式使用美元符号括起，行内公式示例：质能方程：$E=mc^2$
	\item 公式单独显示操作如下：
	
	
	自动编号
	\begin{equation}
		E=mc^2
	\end{equation}
	
	
	不自动编号
	\begin{equation*}
		E=mc^2
	\end{equation*}
	
	
	
	美元符号
	$$E=mc^2$$
	\item 上标，上标只有单个字符时使用“\^{}”符号，如：$E=mc^2$，否则将上标用大括号括起即可，如：$2^{x^{2}+y}$
	\item 下标，下标与上标相似，单个字符使用“\_{}”字符，如$a_1$，否则用大括号括起即可，如：$a_{100}$
	\item 分式：$frac{3}{4}$
	\item 根号：$\sqrt{2}$
	\item 矩阵： 矩阵使用matrix环境排版
	$$\begin{matrix}
		0 & 1 \\
		1 & 0
	\end{matrix}$$
	
	
	patrix环境用于在矩阵两端加小括号
	$$\begin{pmatrix}
		0 & 1 \\
		1 & 0
	\end{pmatrix}$$
	
	
	batrix环境用于在矩阵两端加中括号
	\begin{equation*}
		\begin{bmatrix}
			0 & 1 \\
			1 & 0
		\end{bmatrix}
	\end{equation*}
	
	Batrix环境用于在矩阵两端加大括号
	\begin{equation*}
		\begin{Bmatrix}
			0 & 1 \\
			1 & 0
		\end{Bmatrix}
	\end{equation*}
	
	vatrix环境用于在矩阵两端加竖线
	\begin{equation*}
		\begin{vmatrix}
			0 & 1 \\
			1 & 0
		\end{vmatrix}
	\end{equation*}
	
	Vatrix环境用于在矩阵两端加双竖线
	\begin{equation*}
		\begin{Vmatrix}
			0 & 1 \\
			1 & 0
		\end{Vmatrix}
	\end{equation*}
	
	矩阵中的省略号可用$\ddots,\vdots,\dots$实现
	\iffalse
	矩阵
	&为列分隔符，\\为行分隔符
	\fi
	\item 多行公式，多行公式用gather环境实现
	\begin{gather}
		E=mc^2\\
		E=mc^2
	\end{gather}
	不带编号：
	\begin{gather*}
		E=mc^2\\
		E=mc^2
	\end{gather*}
	单行不编号：使用notag命令
	\begin{gather}
		E=mc^2 \notag \\
		E=mc^2
	\end{gather}
	一个公式的多行编写：
	\begin{equation}
		\begin{split}
			E & =mc^2  \\
			& =mc^2
		\end{split}
	\end{equation}
	\iffalse
	一个公式的多行编写:&为对齐位置
	\fi
	\item 分段函数：
	\begin{equation}
		D(x)=\begin{cases}
			a, &0 \\
			b, &1
			%&为对齐位置
		\end{cases}
	\end{equation} 
\end{enumerate}

\subsubsection{数学公式}

数学公式可以使用 ${equation}$ 和 ${equation*}$ 环境。
注意数学公式的引用应前后带括号，通常使用 ${eqref}$ 命令，比如式\ref{eq:example}。
\iffalse
eqref需要使用标签方式
\fi
\begin{equation}
	\frac{1}{2 \uppi i} \int_\gamma f = \sum_{k=1}^m n(\gamma; a_k) R(f; a_k).
	\label{eq:example}
\end{equation}

多行公式尽可能在“=”处对齐，推荐使用 ${align}$ 环境。
\begin{align}
	a & = b + c + d + e \\
	& = f + g
\end{align}



\subsubsection{数学定理}

定理环境的格式可以使用 ${amsthm}$ 或者 ${ntheorem}$ 宏包配置。
用户在导言区载入这两者之一后，模板会自动配置 ${thoerem}$、${proof}$ 等环境。

\begin{theorem}[Lindeberg--Lévy 中心极限定理]
	设随机变量 $X_1,X_2, \dots, X_n$ 独立同分布， 且具有期望 $\mu$ 和有限的方差 $\sigma^2 \ne 0$，
	记 $\bar{X}_n = \frac{1}{n} \sum_{i+1}^n X_i$，则
	\begin{equation}
		\lim_{n \to \infty} P \left(\frac{\sqrt{n} \left( \bar{X}_n - \mu \right)}{\sigma} \le z \right) = \Phi(z),
	\end{equation}
	其中 $\Phi(z)$ 是标准正态分布的分布函数。
\end{theorem}
\begin{proof}
	Trivial.
\end{proof}

	\chapter{表格与算法}

\section{表格}
自动生成LaTeX表工具: \url{https://www.tablesgenerator.com/}
\par
表应具有自明性。为使表格简洁易读，尽可能采用三线表，如表~\ref{tab:three-line}。
三条线可以使用 booktabs 宏包提供的命令生成。

\begin{table}[htbp]
	\centering
	\caption{三线表示例}
	\begin{tabular}{ll}
		\toprule
		文件名          & 描述                         \\
		\midrule
		zufereport.dtx   & 模板的源文件，包括文档和注释 \\
		zufereport.cls   & 模板文件                     \\
		zufereport-*.bst & BibTeX 参考文献表样式文件    \\
		\bottomrule
	\end{tabular}
	\label{tab:three-line}
\end{table}

表格如果有附注，尤其是需要在表格中进行标注时，可以使用 threeparttable 宏包。
研究生要求使用英文小写字母 a、b、c……顺序编号，本科生使用圈码 ①、②、③……编号。

\begin{table}[htbp]
	\centering
	\begin{threeparttable}[c]
		\caption{带附注的表格示例}
		\label{tab:three-part-table}
		\begin{tabular}{ll}
			\toprule
			文件名                 & 描述                         \\
			\midrule
			zufereport.dtx\tnote{a} & 模板的源文件，包括文档和注释 \\
			zufereport.cls\tnote{b} & 模板文件                     \\
			zufereport-*.bst        & BibTeX 参考文献表样式文件    \\
			\bottomrule
		\end{tabular}
		\begin{tablenotes}
			\item [a] 可以通过 xelatex 编译生成模板的使用说明文档；
			使用 xetex 编译 \file{zufereport.ins} 时则会从 \file{.dtx} 中去除掉文档和注释，得到精简的 \file{.cls} 文件。
			\item [b] 更新模板时，一定要记得编译生成 \file{.cls} 文件，否则编译论文时载入的依然是旧版的模板。
		\end{tablenotes}
	\end{threeparttable}
\end{table}

如某个表需要转页接排，可以使用 \usepackage{longtable} 宏包，需要在随后的各页上重复表的编号。
编号后跟表题（可省略）和“（续）”，置于表上方。续表均应重复表头。

\singlespacing
\begin{longtable}{cccc}
	\caption{跨页长表格的表题}
	\label{tab:longtable} \\
	\toprule
	表头 1 & 表头 2 & 表头 3 & 表头 4 \\
	\midrule
	\endfirsthead
	\caption*{续表~\thetable\quad 跨页长表格的表题} \\
	\toprule
	表头 1 & 表头 2 & 表头 3 & 表头 4 \\
	\midrule
	\endhead
	\bottomrule
	\endfoot
	Row 1  & & & \\
	Row 2  & & & \\
	Row 3  & & & \\
	Row 4  & & & \\
	Row 5  & & & \\
	Row 6  & & & \\
	Row 7  & & & \\
	Row 8  & & & \\
	Row 9  & & & \\
	Row 10 & & & \\
\end{longtable}


\section{伪代码}


算法环境可以使用 algorithms 或者 algorithm2e 宏包。

\renewcommand{\algorithmicrequire}{\textbf{Input：}\unskip}
\renewcommand{\algorithmicensure}{\textbf{Output：}\unskip}
\begin{breakablealgorithm}
	\caption{QuickSort 示例}
	\begin{algorithmic}[1] % 每行显示行号
		\Require $Array$数组，$n$数组大小
		\Ensure 排序后的数组
		\Function{QuickSort}{$Array, left, right$}
		\If {$left < right$}
		\State $pivot \gets$ \Call{Partition}{$Array, left, right$}
		\State \Call{QuickSort}{$Array, left, pivot-1$}
		\State \Call{QuickSort}{$Array, pivot+1, right$}
		\EndIf
		\EndFunction
		\State
		\Function{Partition}{$Array, left, right$}
		\State $pivot \gets Array[right]$
		\State $i \gets left-1$
		\For{$j \gets left$ \textbf{to} $right-1$}
		\If{$Array[j] < pivot$}
		\State $i \gets i+1$
		\State \Call{Swap}{$Array[i], Array[j]$}
		\EndIf
		\EndFor
		\State \Call{Swap}{$Array[i+1], Array[right]}$
		\State \Return{$i+1$}
		\EndFunction
		\State
		\Function{Swap}{$a, b$}
		\State $temp \gets a$
		\State $a \gets b$
		\State $b \gets temp$
		\EndFunction
	\end{algorithmic}

\end{breakablealgorithm}

\section{代码示例}
\lstinputlisting[style=C++, caption={示例 C++ 代码}, label={lst:cpp}]{Codes/test2.cpp}
\lstinputlisting[style=Python,caption={示例 Python 代码},label={lst:py}]{Codes/decision_tree.py}
\lstinputlisting[style=Java,caption={示例 Java 代码},label={lst:java}]{Codes/test.java}

	
	%参考文献
	\newpage
	\printbibliography[title={参考文献}]
	
	%附录
	\newpage
	\begin{appendix}
	\chapter*{Appendix A}
	\addcontentsline{toc}{chapter}{Appendix A}
	这是附录A的内容。
	
	\chapter*{Appendix B}
		\addcontentsline{toc}{chapter}{Appendix B}
	这是附录B的内容。
\end{appendix}

	
	%致谢
	\newpage
	\chapter*{Acknowledgements}
\addcontentsline{toc}{chapter}{Acknowledgements}

\lipsum[1]
\end{document}
